\chapter{Instrukcje warunkowe}

\section{Po co programowi warunki}
Strona internetowa rzadko jest absolutnie identyczna w kazdej sytuacji. Czasem chcemy:

\begin{itemize}
  \item pokazac sekcje tylko dla opublikowanych wpisow,
  \item ukryc komunikat, jesli lista jest pusta,
  \item wyswietlic inny tekst dla uzytkownika premium,
  \item nadac inny styl w zaleznosci od stanu danych.
\end{itemize}

Do takich sytuacji sluzy instrukcja \mc{if}. Pozwala powiedziec:
\emph{jesli warunek jest prawdziwy, wykonaj ten blok; w przeciwnym razie wykonaj inny
albo nie wykonuj niczego}.

\section{Najprostsza forma \mc{if}}
\begin{lstlisting}[style=moonchunk,caption={Podstawowy warunek if}]
if (isPublished) {
  print("Wpis jest opublikowany");
};
\end{lstlisting}

Jesli \mc{isPublished} ma wartosc \mc{true}, instrukcja wewnatrz bloku zostanie wykonana.
Jesli ma wartosc \mc{false}, nic sie nie stanie.

\section{Warunek musi byc typu \mc{bool}}
To jedna z bardzo waznych zasad MoonChunk. Wyrazenie w nawiasach po \mc{if} powinno
po obliczeniu dawac \mc{true} albo \mc{false}. Nie nalezy traktowac liczb albo tekstow
jako niejawnych zastepnikow prawdy.

Poprawnie:

\begin{lstlisting}[style=moonchunk,caption={Poprawny warunek logiczny}]
let showBanner: bool = views > 100;
if (showBanner) {
  print("Pokaz baner");
};
\end{lstlisting}

Niepoprawnie:

\begin{lstlisting}[style=moonchunk,caption={Niepoprawny warunek niebedacy bool}]
if (3) {
  print("To powinno byc bledem");
};
\end{lstlisting}

\section{\mc{if} z \mc{else}}
Bardzo czesto potrzebujemy dwoch alternatyw. Wtedy korzystamy z \mc{else}.

\begin{lstlisting}[style=moonchunk,caption={Pelny warunek if-else}]
if (isPublished) {
  print("Status: opublikowany");
} else {
  print("Status: szkic");
};
\end{lstlisting}

To prosty, bardzo czytelny zapis: jedna droga dla prawdy, druga dla falszu.

\section{Warunki a HTML}
Instrukcje warunkowe sa szczegolnie przydatne przy generowaniu tresci strony.
Mozesz decydowac, ktory fragment \mc{content} ma zostac dolaczony.

\begin{lstlisting}[style=moonchunk,caption={Warunkowe generowanie fragmentu strony}]
page "" {
  let hasPosts: bool = true;

  if (hasPosts) {
    content {
      <p>Mamy wpisy do pokazania.</p>
    };
  } else {
    content {
      <p>Brak wpisow.</p>
    };
  };
};
\end{lstlisting}

To jeden z najwazniejszych praktycznych wzorcow w MoonChunk.

\section{Warunki zlozone}
Prawdziwa sila \mc{if} pojawia sie wtedy, gdy warunek sklada sie z kilku elementow.
Na przyklad:

\begin{lstlisting}[style=moonchunk,caption={Zlozony warunek logiczny}]
if (isPublished and not isArchived and views > 100) {
  print("Pokaz promowany wpis");
};
\end{lstlisting}

Takie warunki sa bardzo przydatne, ale trzeba pilnowac czytelnosci. Jesli logika staje
sie dluga, warto podzielic ja na zmienne pomocnicze.

\begin{lstlisting}[style=moonchunk,caption={Podzial zlozonego warunku na kroki}]
let canPromote: bool = isPublished and not isArchived;
let isPopular: bool = views > 100;

if (canPromote and isPopular) {
  print("Pokaz promowany wpis");
};
\end{lstlisting}

\section{Zagniezdzone warunki}
Mozna tez umieszczac jeden \mc{if} wewnatrz drugiego. Trzeba jednak robic to z umiarem.

\begin{lstlisting}[style=moonchunk,caption={Zagniezdzone warunki}]
if (isPublished) {
  if (views > 100) {
    print("Wpis jest opublikowany i popularny");
  };
};
\end{lstlisting}

Zagniezdzenia sa czasem naturalne, ale ich naduzywanie pogarsza czytelnosc. Często
lepszym rozwiazaniem jest obliczenie kilku nazwanych warunkow pomocniczych.

\section{If a zakresy nazw}
W MoonChunk ciało instrukcji \mc{if} jest miejscem, w ktorym trzeba uwazac na redeklaracje
nazw. Jezyk celowo nie zachowuje sie tu jak czysty, niezalezny blok do bezkarnego
zaslaniania tych samych nazw z gory.

Przyklad potencjalnie problematyczny:

\begin{lstlisting}[style=moonchunk,caption={Redeclaracja w ciele if}]
const x: int = 3;

if (true) {
  const x: int = 5;
};
\end{lstlisting}

Taki zapis nalezy traktowac jako blad redeklaracji. Jesli chcesz utworzyc odrebny,
celowo zagniezdzony zakres, skorzystaj z jawnego bloku \mc{\{ ... \}}, o czym bedzie
mowa w osobnym rozdziale.

\section{Kiedy uzywac \mc{if}, a kiedy operatora \mc{?:}}
To czeste pytanie. Dobra zasada jest bardzo praktyczna:

\begin{itemize}
  \item operator \mc{?:} jest dobry dla bardzo prostego wyboru jednej z dwoch wartosci,
  \item \mc{if} jest lepszy, gdy trzeba wykonac wiecej niz jedna instrukcje albo zbudowac HTML.
\end{itemize}

Przyklad prosty:

\begin{lstlisting}[style=moonchunk,caption={Operator warunkowy do malego wyboru}]
let label: string = isPublished ? "Tak" : "Nie";
\end{lstlisting}

Przyklad bardziej rozbudowany:

\begin{lstlisting}[style=moonchunk,caption={if do wiekszej logiki}]
if (isPublished) {
  print("Wpis gotowy");
  content {
    <span>Opublikowany</span>
  };
} else {
  print("Wpis jeszcze nie gotowy");
};
\end{lstlisting}

\section{Typowe zastosowania w stronach}
W prawdziwych projektach \mc{if} jest szczegolnie przydatne do:

\begin{itemize}
  \item pokazywania pustego stanu listy,
  \item warunkowego renderowania banerow,
  \item wyboru sekcji dla roznych typow tresci,
  \item nadpisywania wybranych fragmentow na podstawie danych,
  \item ochrony przed probą dostepu do danych, ktorych moze nie byc.
\end{itemize}

\section{Najczestsze bledy przy warunkach}
\begin{itemize}
  \item wpisanie \mc{=} zamiast \mc{==},
  \item uzycie liczby lub tekstu zamiast wartosci logicznej,
  \item zbyt rozbudowany warunek bez nazw pomocniczych,
  \item redeklarowanie tej samej nazwy w ciele \mc{if},
  \item zapominanie o srednikach po blokach, jesli skladnia ich wymaga.
\end{itemize}

Jesli MoonChunk zglosi blad w warunku, warto najpierw sprawdzic trzy rzeczy:

\begin{enumerate}
  \item czy porownanie jest zapisane jako \mc{==}, a nie \mc{=},
  \item czy wynik warunku naprawde jest typu \mc{bool},
  \item czy nawiasy i klamry zamykaja sie poprawnie.
\end{enumerate}
